\section{De Demo a producción: los humanos siguen a cargo de la Objective Function}

\subsection{El trabajo humano no desaparece, solo se desplaza a la capa de definición de objetivos}
Karpathy insiste en varias ocasiones en que se debe retirar al humano del loop, pero esto no significa que los humanos dejen de ser importantes. Al contrario, la responsabilidad humana se desplaza de ``ejecutar cada paso'' a ``definir el objective, establecer los límites y proporcionar el verifier''. Auto research funciona no porque el agent produzca valor de manera mágica, sino porque la función objetivo es suficientemente clara.

\begin{importantbox}{Lo que más teme un agent no es la falta de capacidad, sino la ambigüedad de la función objetivo}
\begin{itemize}
    \item Escribir un CUDA kernel, pasar el benchmark más rápido, reducir el loss de entrenamiento: este tipo de tareas se puede juzgar con métricas objetivas;
    \item el gusto del producto, la intención del usuario y las prioridades de la organización, en cambio, muchas veces no tienen una única respuesta correcta;
    \item en cuanto la objective function no es clara, el sistema termina desperdiciando una gran cantidad de tokens en óptimos locales y direcciones equivocadas.
\end{itemize}
\end{importantbox}

Esta es también la razón por la que él pone juntas las frases ``I shouldn't be looking at the results'' y ``if you can't evaluate, you can't auto research it''. La primera habla de eliminar el cuello de botella humano, y la segunda habla de los límites de aplicabilidad de la automatización. Combinadas, forman un juicio muy propio de la ingeniería: \textbf{todo lo que se pueda verificar debe, en la medida de lo posible, incorporarse al ciclo cerrado; todo lo que no se pueda verificar no debe fingirse como ya automatizado.}

\subsection{De la API agent-first a la infraestructura de la empresa}
Si en el futuro ``apps shouldn't exist'', entonces el activo más importante de una empresa de software ya no será solo su interfaz frontal, sino la API, el modelo de permisos de datos, la observabilidad y los registros de auditoría que hay detrás de ella. Para que el agent se convierta en el usuario principal, el sistema debe exponerse de forma nativa al agent, en lugar de depender de la automatización de la UI para operar de manera inversa.

\begin{knowledgebox}{Lista de infraestructura para productos agent-first}
\begin{enumerate}
    \item \textbf{API componible}: estable, autorizable y auditable;
    \item \textbf{Capa de estado y memoria}: permite que el agent continúe tareas entre sesiones;
    \item \textbf{Observabilidad}: saber qué hizo el agent y por qué falló;
    \item \textbf{Reversión y control de permisos}: evita que la automatización de cadenas largas amplifique los errores.
\end{enumerate}
\end{knowledgebox}

\begin{table}[H]
\centering
\begin{tabular}{p{3.0cm}p{4.1cm}p{4.8cm}}
\toprule
\textbf{Pregunta} & \textbf{Supuesto del software antiguo} & \textbf{Supuesto agent-first} \\
\midrule
¿Quién es el usuario principal? & El humano hace clic en la UI & El agent llama a la API, el humano supervisa y acepta \\
¿Cuál es el punto de entrada del producto? & Páginas, botones, formularios & message channel, API, workflow trigger \\
¿Cómo se mejora la eficiencia? & Simplificando el flujo de interacción & Exponiendo capacidades operativas, reduciendo el cambio de contexto \\
¿Cómo se acumulan barreras de entrada? & Una interfaz frontal y procesos más refinados & Datos, integración de herramientas, modelo de permisos y evaluation loop \\
\bottomrule
\end{tabular}
\caption{Lo que Karpathy llama agent-first no es simplemente agregar una ventana de chat más, sino reescribir la lógica de distribución del software}
\end{table}

\subsection{Resumen de la sección}
Los humanos no desaparecerán del sistema, pero se ubicarán más en la posición de la función objetivo, los límites de permisos y los criterios de aceptación. El verdadero centro de la puesta en producción consiste en conectar las capacidades del agent a la API, la gestión de estado y la observabilidad, en lugar de quedarse en una automatización de tipo Demo.


